% ---------------- Title page ----------------
\begin{titlepage}
\centering
{\large TERM THESIS REPORT\par}
\vspace{1.2cm}
{\LARGE\bfseries \thesistitle\par}
\vspace{1.4cm}
{\large Submitted in partial fulfilment of the requirements for the degree of\par}
\vspace{0.3cm}
{\large\bfseries \degree\par}
\vspace{1.2cm}
{\large by\par}
\vspace{0.3cm}
{\Large\bfseries \authorname\par}
{\large (\rollno)\par}
\vspace{1.2cm}
{\large Under the supervision of\par}
\vspace{0.3cm}
{\Large \supervisor\par}
\vfill
% \includegraphics[width=3cm]{iitkgp_logo}  % add logo file to figures/ if required
{\large \department\par}
{\large\bfseries Indian Institute of Technology Kharagpur\par}
{\large Kharagpur -- 721302, West Bengal, India\par}
\vspace{0.4cm}
{\large \submitdate\par}
\end{titlepage}

% ---------------- Certificate ----------------
\chapter*{Certificate}
\addcontentsline{toc}{chapter}{Certificate}
This is to certify that the term thesis entitled \textbf{``\thesistitle''}, submitted by
\textbf{\authorname} (Roll No.\ \rollno) to the \department, Indian Institute of Technology
Kharagpur, is a record of bona fide work carried out under my supervision and guidance
during the academic session 2026--27.

\vspace{3cm}
\noindent\begin{tabular}{@{}l}
\rule{6cm}{0.4pt}\\
\supervisor\\
\department\\
Indian Institute of Technology Kharagpur\\
Date:
\end{tabular}

% ---------------- Declaration ----------------
\chapter*{Declaration}
\addcontentsline{toc}{chapter}{Declaration}
I certify that (a) the work contained in this report is original and has been done by me under
the guidance of my supervisor; (b) the work has not been submitted to any other institute for
any degree or diploma; (c) I have followed the guidelines provided by the Institute in preparing
the report; (d) whenever I have used materials (data, theoretical analysis, figures and text)
from other sources, I have given due credit to them by citing them in the text of the report
and giving their details in the references; and (e) all data are from openly licensed public
sources, and all code was written for this project and is available for verification.

\vspace{3cm}
\noindent\begin{tabular}{@{}l}
\rule{6cm}{0.4pt}\\
\authorname\\
\rollno
\end{tabular}

% ---------------- Acknowledgements ----------------
\chapter*{Acknowledgements}
\addcontentsline{toc}{chapter}{Acknowledgements}
I express my sincere gratitude to my supervisor, \supervisor, for guidance and encouragement
throughout this work. I thank ADEME and INRAE for making the AGRIBALYSE database openly
available, and the developers of the open-source scientific Python ecosystem
(scikit-learn, XGBoost, SHAP, Sentence-Transformers, Streamlit) on which this work is built.
% Personalise as appropriate.

% ---------------- Abstract ----------------
\chapter*{Abstract}
\addcontentsline{toc}{chapter}{Abstract}
Life Cycle Assessment (LCA) is the standard method for quantifying the environmental impacts
of a product from raw-material extraction to end of life, but a full ISO 14040/44 study is
slow, data-hungry and expensive. In practice, designers and analysts need fast
\emph{screening} estimates long before complete inventory data exist, and must manually map
thousands of product descriptions to database processes. This thesis investigates how
Artificial Intelligence and Machine Learning (AI/ML) can be integrated into LCA to address
these two bottlenecks, and builds a working, openly reproducible prototype.

Using the open AGRIBALYSE~3.1 database (2{,}451 food products, 16 Environmental Footprint
indicators), two modules are developed. \textbf{Module~A} is a surrogate impact model: a
gradient-boosted tree regressor (XGBoost) that predicts a product's impact using only
information available \emph{before} an LCA is done (food category, storage mode,
consumer preparation, recipe ingredients and words in the product name), with a strict
no-leakage design. \textbf{Module~B} is a semantic mapping engine that uses multilingual
sentence embeddings (E5) to match free-text descriptions, in English or French, to the
closest LCA datasets. Predictions are made explainable with SHAP and given calibrated
90\% uncertainty intervals with Mondrian conformal prediction.

Against the common practitioner shortcut of borrowing the sub-category average, the
model raises \Rsq{} (log scale) for climate change from 0.62 to 0.84 and the share of
predictions within a factor of two from 77\% to 91\% under random cross-validation; for
food sub-categories never seen in training, \Rsq{} rises from 0.33 to 0.62. Gains hold for
all six impact categories tested, including water use (\Rsq{} 0.27~$\rightarrow$~0.64). The
conformal intervals achieve 91\% empirical coverage at the 90\% target. The mapping engine
places the correct dataset in its top five for 95\% of terse three-word queries, and five
encoders are benchmarked. SHAP recovers known LCA patterns (ruminant meat raises, plant
ingredients lower the footprint). Both modules are integrated into an interactive
screening tool that flags inputs outside its applicability domain. The thesis closes with a
critical discussion of limitations, including geographic transferability to India, and a
roadmap for future work.

\vspace{0.5cm}
\noindent\textbf{Keywords:} Life Cycle Assessment; machine learning; surrogate modelling;
explainable AI; conformal prediction; sentence embeddings; AGRIBALYSE; food carbon footprint.
